% GENERATED FILE. Edit canonical lessons, then run npm run generate:books.
% canonical-source-hash: fnv1a64:b1e083c4c18d3d17
% canonical-lessons: RU-C198-passazhir, RU-C198-vagon, RU-C198-peresadka, RU-C198-probka, RU-C198-avariya

\chapter{Passenger, Carriage, Transfer, Traffic jam, Accident}
\label{ch:ru-passenger-carriage-transfer-traffic-jam-accident}
\hlchaptermodality{198}

\begin{chapteropening}
\textbf{By the end of this chapter:} \emph{I can say a passenger, a carriage, a transfer, a traffic jam and an accident.}

Complete the last lesson of chapter 198: I can say a passenger, a carriage, a transfer, a traffic jam and an accident.
\end{chapteropening}

\section[\texorpdfstring{\ru{пассажир} (passazhír)}{пассажир (passazhír)}]{\ru{пассажир} (passazhír) — a passenger}

\label{lesson:RU-C198-passazhir}



\begin{quote}
Before the new one: say the Russian for a guitar, then the Russian for a piano.
\end{quote}

\subsection*{You'll want to know: \ru{пассажир}}
\textbf{\ru{пассажир}} (\emph{passazhír}) — \textquotedblleft{}a passenger\textquotedblright{}.

\begin{cousinweb}
\textbf{\ru{пассажир}} — first \emph{\ru{пасажера}}, in 1706; it came through German \emph{Passagier} or Dutch \emph{passagier} from French \emph{passager}, Italian \emph{passaggiere}, \textquotedblleft{}one passing through\textquotedblright{} (Vasmer).
\end{cousinweb}

\begin{grammarlens}[title={What you've built}]
One more word for this chapter.
\end{grammarlens}

\subsection*{Your turn}
\begin{itemize}
\raggedright
  \item \emph{Say it:} \emph{passazhír}
  \item \emph{Say it:} \emph{passazhír}, once more
  \item \emph{Say it:} say \emph{passazhír}
  \item \emph{Recall:} say the Russian for a guitar, then the Russian for a piano, then say \emph{passazhír} again
\end{itemize}

\begin{tcolorbox}[breakable,colback=teal!4,colframe=teal!35!black,title={Before you move on}]
What does \ru{пассажир} mean? (\textquotedblleft{}A passenger\textquotedblright{}.) Say it once more.
\end{tcolorbox}
\section[\texorpdfstring{\ru{вагон} (vagón)}{вагон (vagón)}]{\ru{вагон} (vagón) — a carriage (of a train)}

\label{lesson:RU-C198-vagon}



\begin{quote}
Before the new one: say the Russian for a piano, then the Russian for a passenger.
\end{quote}

\subsection*{You'll want to know: \ru{вагон}}
\textbf{\ru{вагон}} (\emph{vagón}) — \textquotedblleft{}a carriage (of a train)\textquotedblright{}.

\begin{cousinweb}
\textbf{\ru{вагон}} came from German \emph{Waggon} or French \emph{wagon}, which go back to English \emph{waggon}; German \emph{Wagen} and Russian \emph{\ru{воз}}, \textquotedblleft{}a cart\textquotedblright{}, are cousins of the English word (Vasmer).
\end{cousinweb}

\begin{grammarlens}[title={What you've built}]
One more word for this chapter.
\end{grammarlens}

\subsection*{Your turn}
\begin{itemize}
\raggedright
  \item \emph{Say it:} \emph{vagón}
  \item \emph{Say it:} \emph{vagón}, once more
  \item \emph{Say it:} say \emph{vagón}
  \item \emph{Recall:} say the Russian for a piano, then the Russian for a passenger, then say \emph{vagón} again
\end{itemize}

\begin{tcolorbox}[breakable,colback=teal!4,colframe=teal!35!black,title={Before you move on}]
What does \ru{вагон} mean? (\textquotedblleft{}A carriage (of a train)\textquotedblright{}.) Say it once more.
\end{tcolorbox}
\section[\texorpdfstring{\ru{пересадка} (peresádka)}{пересадка (peresádka)}]{\ru{пересадка} (peresádka) — a transfer, a change (of trains)}

\label{lesson:RU-C198-peresadka}



\begin{quote}
Before the new one: say the Russian for a passenger, then the Russian for a carriage.
\end{quote}

\subsection*{You'll want to know: \ru{пересадка}}
\textbf{\ru{пересадка}} (\emph{peresádka}) — \textquotedblleft{}a transfer, a change (of trains)\textquotedblright{}.

\begin{cousinweb}
\textbf{\ru{пересадка}} is \emph{\ru{пере}-}, \textquotedblleft{}across\textquotedblright{}, with \emph{-\ru{садка}}, from \emph{\ru{садить}}, \textquotedblleft{}to seat\textquotedblright{}: changing from one seat to another.
\end{cousinweb}

\begin{grammarlens}[title={What you've built}]
One more word for this chapter.
\end{grammarlens}

\subsection*{Your turn}
\begin{itemize}
\raggedright
  \item \emph{Say it:} \emph{peresádka}
  \item \emph{Say it:} \emph{peresádka}, once more
  \item \emph{Say it:} say \emph{peresádka}
  \item \emph{Recall:} say the Russian for a passenger, then the Russian for a carriage, then say \emph{peresádka} again
\end{itemize}

\begin{tcolorbox}[breakable,colback=teal!4,colframe=teal!35!black,title={Before you move on}]
What does \ru{пересадка} mean? (\textquotedblleft{}A transfer, a change (of trains)\textquotedblright{}.) Say it once more.
\end{tcolorbox}
\section[\texorpdfstring{\ru{пробка} (próbka)}{пробка (próbka)}]{\ru{пробка} (próbka) — a traffic jam; a cork}

\label{lesson:RU-C198-probka}



\begin{quote}
Before the new one: say the Russian for a carriage, then the Russian for a transfer.
\end{quote}

\subsection*{You'll want to know: \ru{пробка}}
\textbf{\ru{пробка}} (\emph{próbka}) — \textquotedblleft{}a traffic jam; a cork\textquotedblright{}.

\begin{cousinweb}
\textbf{\ru{пробка}} apparently came through Low German \emph{propke}, \emph{prop}, \textquotedblleft{}a stopper\textquotedblright{}, since Russia long imported corks from Bremen (Vasmer). A traffic jam corks the road.
\end{cousinweb}

\begin{grammarlens}[title={What you've built}]
One more word for this chapter.
\end{grammarlens}

\subsection*{Your turn}
\begin{itemize}
\raggedright
  \item \emph{Say it:} \emph{próbka}
  \item \emph{Say it:} \emph{próbka}, once more
  \item \emph{Say it:} say \emph{próbka}
  \item \emph{Recall:} say the Russian for a carriage, then the Russian for a transfer, then say \emph{próbka} again
\end{itemize}

\begin{tcolorbox}[breakable,colback=teal!4,colframe=teal!35!black,title={Before you move on}]
What does \ru{пробка} mean? (\textquotedblleft{}A traffic jam; a cork\textquotedblright{}.) Say it once more.
\end{tcolorbox}
\section[\texorpdfstring{\ru{авария} (aváriya)}{авария (aváriya)}]{\ru{авария} (aváriya) — an accident, a crash}

\label{lesson:RU-C198-avariya}



\begin{quote}
Before the new one: say the Russian for a transfer, then the Russian for a traffic jam; a cork.
\end{quote}

\subsection*{You'll want to know: \ru{авария}}
\textbf{\ru{авария}} (\emph{aváriya}) — \textquotedblleft{}an accident, a crash\textquotedblright{}.

\begin{cousinweb}
\textbf{\ru{авария}} comes from Italian \emph{avaria} or French \emph{avarie}, \textquotedblleft{}damage to a ship or cargo\textquotedblright{}, from Arabic \emph{ʿawār}, \textquotedblleft{}damage, a flaw\textquotedblright{} (Vasmer).
\end{cousinweb}

\begin{grammarlens}[title={What you've built}]
One more word for this chapter.
\end{grammarlens}

\subsection*{Your turn}
\begin{itemize}
\raggedright
  \item \emph{Say it:} \emph{aváriya}
  \item \emph{Say it:} \emph{aváriya}, once more
  \item \emph{Say it:} say \emph{aváriya}
  \item \emph{Recall:} say the Russian for a transfer, then the Russian for a traffic jam; a cork, then say \emph{aváriya} again
\end{itemize}

\begin{tcolorbox}[breakable,colback=teal!4,colframe=teal!35!black,title={Before you move on}]
What does \ru{авария} mean? (\textquotedblleft{}An accident, a crash\textquotedblright{}.) Say it once more.
\end{tcolorbox}
